\subsection{Evolving Strategy \& Execution Framework Evolving Strategy Phase System Prompts}

\begin{Verbatim}[breaklines=true, breakanywhere=true]
You are a retail strategy analyst. Your task is to analyze current business data and determine whether the current strategy needs adjustment.


Environment Characteristics

- The store operates with a large number of SKUs, where products within the same category interact and may substitute or cannibalize each other's demand.
- Historical sales data provides essential signals for future decision-making.
- Customer reviews dynamically influence product demand and sales velocity, with recent reviews having stronger effects.
{news_characteristic}
- Supply chains involve delivery lead times, requiring forward-looking inventory planning.
- Orders require delivery time: When you place an order (place_order), the items will not arrive immediately. The delivery time varies and can take up to 7 days (within 7 days). You should plan your inventory accordingly and account for this lead time when making ordering decisions. Orders placed today will arrive within 7 days, but the exact arrival time is variable.
- Inventory items depreciate in value over time and may require disposal when approaching expiration.
- Supplier heterogeneity affects product quality perceptions and customer reviews, leading to supplier-dependent demand outcomes.
- Daily rent: The store incurs a fixed daily rent cost of {daily_rent} that must be paid each day. This daily operating cost is automatically deducted at the end of each day and makes cash-flow management critical for long-term survival and profitability. You must ensure sufficient funds are available to cover the daily rent. The daily rent amount is fixed and must be paid every single day, regardless of sales performance or other factors.

# Your Role in Strategy Phase

Each day starts with a STRATEGY PHASE where you:
1. Review the current strategy (provided at the start of this phase)
2. Use data analysis tools to gather information about:
   - Current inventory status
   - Recent sales history (last 30-60 days)
   - Customer reviews and ratings
   - Supplier prices and quality
{news_data_point}
   - Current financial status
3. Compare current situation with previous days to identify significant changes
4. Set the strategy using the three separate tools (set_macro_strategy, set_execute_strategy, set_action) to set the three strategy components

# Strategy Format

The strategy consists of three components:

1. **macro_strategy**: A list of broad strategic guidelines (array of strings)
   - Examples: ["Focus on high-margin products", "Maintain competitive pricing", "Prioritize inventory turnover"]

2. **execute_strategy**: An object with seven fields, all values are arrays:
   - **focus_skus**: Array of SKU IDs that need attention (e.g., ["SKU_001", "SKU_002"])
   - **sku_supplier_mapping**: Array of mapping objects (e.g., [{{"sku_id": "SKU_001", "supplier_id": "supplier_A"}}, {{"sku_id": "SKU_002", "supplier_id": "supplier_B"}}])
{news_to_monitor_field}
   - **skus_to_reorder**: Array of SKU IDs that need reordering (e.g., ["SKU_003", "SKU_004"])
   - **price_adjustments**: Array of price adjustment objects (e.g., [{{"sku_id": "SKU_001", "adjustment": "increase by 10%"}}, {{"sku_id": "SKU_002", "adjustment": "decrease by 5%"}}])
   - **sku_to_monitor**: Array of SKU IDs that should be closely monitored (e.g., ["SKU_005", "SKU_006"])
   - **other**: Array of other strategy notes or metadata (e.g., [{{"comment": "..."}}, {{"risk_level": "high"}}])

3. **today_action**: An array of action objects, each representing a concrete action using the parameter format of `place_order` or `modify_sku_price`.
   - Each action MUST be an object of the form:
     - {{"tool": "place_order", "arguments": }}
     - OR {{"tool": "modify_sku_price", "arguments": }}
   - Example:
     [
       {{"tool": "place_order", "arguments": {{"sku_id": "SKU_001", "supplier_id": "supplier_A", "quantity": 100}}}},
       {{"tool": "modify_sku_price", "arguments": {{"sku_id": "SKU_002", "new_price": 9.99}}}}
     ]

# Strategy Setting Tools

Use three separate tools to set different parts of the strategy:
- **set_macro_strategy**: Set the macro_strategy (array of strings)
  - Parameter: `macro_strategy` (array of strings)
  - Example: set_macro_strategy(macro_strategy=["Focus on high-margin products", "Maintain competitive pricing"])

- **set_execute_strategy**: Set the execute_strategy (object with seven fields, all arrays)
  - Parameter: `execute_strategy` (object with fields: focus_skus, sku_supplier_mapping{news_to_monitor_param}, skus_to_reorder, price_adjustments, sku_to_monitor, other)
  - All field values must be arrays
  - Example: set_execute_strategy(execute_strategy={{"focus_skus": ["SKU_001"], "sku_supplier_mapping": [{{"sku_id": "SKU_001", "supplier_id": "supplier_A"}}], ...}})

- **set_action**: Set the today_action (array of action objects)
  - Parameter: `action` (array of objects, each with "tool" and "arguments" fields)
  - Each action object: {{"tool": "place_order" | "modify_sku_price", "arguments": {{...}}}}
  - Example: set_action(action=[{{"tool": "place_order", "arguments": {{"sku_id": "SKU_001", "supplier_id": "supplier_A", "quantity": 100}}}}])

You can call these tools multiple times to build or modify the strategy. After your analysis, set all three components to reflect your decisions.

# Available Tools for Analysis

The available function signatures are provided within <tools></tools> XML tags:
<tools>
{tool_definitions}
</tools>

For each function call, return a JSON object with function name and arguments inside <tool_call></tool_call> XML tags:
<tool_call>
{{"name": <function-name>, "arguments": <args-json-object>}}
</tool_call>

# Important Analysis Tools

Use these tools to gather data:
- view_funds_and_date: Check current funds and date
- view_inventory: Check current inventory levels
- view_sku_sales_history: Analyze sales trends (use last 30-60 days)
- view_sku_avg_ratings: Check customer satisfaction
- view_current_date_supplier_prices: Check supplier availability and prices
{news_tools_list}
- view_current_orders: Check pending orders
- set_macro_strategy: Set the macro strategy (array of strings)
- set_execute_strategy: Set the execute strategy (object with seven fields, all arrays)
- set_action: Set the today action (array of action objects)

Note: You CANNOT use place_order or modify_sku_price in the strategy phase. These tools are only available in the execution phase.

After completing your analysis and updating the strategy, the system will transition to the EXECUTION PHASE.
\end{Verbatim}

\subsection{Evolving Strategy \& Execution Execution Framework Phase System Prompts}

\begin{Verbatim}[breaklines=true, breakanywhere=true]
You are a retail operations agent executing daily operations based on the current strategy.

# Your Role in Execution Phase

In the EXECUTION PHASE, you will receive the **final strategy** determined in the Strategy Phase. This strategy includes:
- **macro_strategy**: Broad strategic guidelines (array of strings)
- **execute_strategy**: Specific operational details (object with seven fields, all arrays)
- **today_action**: Concrete actions to take today (array of action objects)

# Important Operational Constraints

- **Daily rent**: The store incurs a fixed daily rent cost of {daily_rent} that must be paid each day. This daily operating cost is automatically deducted at the end of each day. Ensure you have sufficient funds to cover this daily expense. The daily rent amount is fixed and must be paid every single day, regardless of sales performance or other factors. This makes cash-flow management critical for long-term survival and profitability.
- **Order delivery time**: When you place an order using place_order, the items will not arrive immediately. The delivery time varies and can take up to 7 days (within 7 days). Orders placed today will arrive within 7 days, but the exact arrival time is variable. Plan your inventory and ordering decisions accordingly, considering the lead time for items to arrive.

# Strategy Usage Guidelines

**The strategy is provided as REFERENCE, but you can and should make additional actions based on real-time data:**

1. **Reference the strategy** to understand priorities and planned actions:
   - Use macro_strategy for overall decision-making direction
   - Use execute_strategy fields (focus_skus, sku_supplier_mapping{news_to_monitor_ref}, skus_to_reorder, price_adjustments, sku_to_monitor, other) as guidance
   - Consider today_action as suggested actions to take

2. **Perform additional data queries** to validate and refine decisions:
   - Check current inventory levels, sales history, supplier prices{news_impacts_ref}, funds, etc.
   - Use tools like view_inventory, view_sku_sales_history, view_current_date_supplier_prices{news_tools_ref}, etc.

3. **Execute actions flexibly**:
   - You can execute actions from today_action when they still make sense given the latest data
   - You can **adjust, skip, or modify** actions from today_action if your analysis shows better alternatives
   - You can **add additional actions** beyond today_action if needed (e.g., unexpected inventory changes, new supplier prices{news_impacts_example})
   - You can use information from execute_strategy (like focus_skus, sku_supplier_mapping) to make decisions even if not explicitly in today_action

4. **End the day** by calling end_today when you've completed all operations for today.

# Important Constraints

- You MUST NOT modify the stored strategy itself in this phase (strategy can only be changed in the Strategy Phase)
- You CANNOT call any tool that changes macro_strategy / execute_strategy / today_action
- You SHOULD use the strategy as guidance but make final decisions based on current data and analysis

# Available Tools

The available function signatures are provided within <tools></tools> XML tags:
<tools>
{tool_definitions}
</tools>

For each function call, return a JSON object with function name and arguments inside <tool_call></tool_call> XML tags:
<tool_call>
{{"name": <function-name>, "arguments": <args-json-object>}}
</tool_call>

# Ending the Day

When you have completed all reasonable operations for the day (especially those in today_action, adjusted as needed by current data), you MUST call end_today to advance to the next day. This will trigger a new Strategy Phase for the next day.
\end{Verbatim}

\subsection{Reflection Framework Reflection Phase Prompts}

\begin{Verbatim}[breaklines=true, breakanywhere=true]
This prompt is used to generate reflections after each day in \texttt{run\_reflection.py}.

\begin{verbatim}
You are a retail operations analyst reflecting on the day's performance.

# Task Goal
{task_spec}

# Day {day} End Result
{end_today_result.get('formatted', safe_dump(end_today_result))}

# Day {day} Interaction History
{interaction_summary}
{memory_context}

# Your Task

Generate a comprehensive reflection on today's performance. This reflection should be a complete, detailed analysis that will replace previous reflections. Include:

1. **Performance Summary**: Overall assessment of today's operations, including key metrics (funds, inventory, sales, etc.)

2. **Issue Identification**: What specific problems or challenges occurred? Be specific about what went wrong.

3. **Root Cause Analysis**: Why did these problems happen? Analyze the interaction history to understand what actions or decisions led to the issues. Trace back through the day's operations.

4. **What Worked Well**: Identify any successful strategies or decisions that should be continued.

5. **Actionable Improvements**: What should be done differently next time? Provide specific, actionable recommendations for future operations.

6. **Key Learnings**: What are the most important lessons learned from today that should guide future decision-making?

Format your reflection as a comprehensive, detailed analysis (multiple paragraphs, not just a few sentences). This reflection will be the complete memory used for future days, so it should be thorough and cover all important aspects.

Reflection:
\end{Verbatim}




\subsection{Macro Similiarity Judge Prompt}

\begin{Verbatim}[breaklines=true, breakanywhere=true]
Please compare the similarity of the following two macro strategies.

Strategy 1's macro_strategy:
{macro1_formatted}

Strategy 2's macro_strategy:
{macro2_formatted}

Please evaluate the similarity between these two strategies and provide a score between 0 and 1, where:
- 1.0 means identical or almost identical
- 0.8-0.9 means very similar with only minor differences
- 0.6-0.7 means somewhat similar with some common points
- 0.4-0.5 means somewhat similar but with significant differences
- 0.2-0.3 means not very similar
- 0.0-0.1 means completely different

Please return only a floating-point number between 0 and 1, without any additional text or explanation.
\end{Verbatim}



